% GENERATED FILE. Edit canonical lessons, then run npm run generate:books.
% canonical-source-hash: fnv1a64:b01e349ae087d574
% canonical-lessons: GU-C259-mec, GU-C259-haravvu, GU-C259-dodvu, GU-C259-kasrat, GU-C259-kusti

\chapter{Matches and wrestling}
\label{ch:gu-matches-and-wrestling}
\hlchaptermodality{259}

\begin{chapteropening}
\textbf{By the end of this chapter:} \emph{I can say who won the match, who ran and what exercise I did.}

Complete the last lesson of chapter 259: I can say who won the match, who ran and what exercise I did.
\end{chapteropening}

\section[\texorpdfstring{\gu{મેચ} (mec)}{મેચ (mec)}]{\gu{મેચ} (mec) — match}

\label{lesson:GU-C259-mec}



\begin{quote}
Before the new one: say the Gujarati for to hit, then the Gujarati for to bounce.
\end{quote}

\subsection*{You'll want to know: \gu{મેચ}}
\textbf{\gu{મેચ}} — \emph{mec} — \textquotedblleft{}match\textquotedblright{}.

\begin{quote}
\textbf{\gu{ગઈકાલે} \gu{અમારી} \gu{શાળાની} \gu{મેચ} \gu{હતી}.} — \emph{gaīkāle amārī śāḷānī mec hatī} — \textquotedblleft{}Yesterday our school had a match.\textquotedblright{}
\end{quote}

\begin{cousinweb}
\textbf{\gu{મેચ}} (\emph{mec}) is a loan from English \emph{match}, a game between two sides.
\end{cousinweb}

\begin{grammarlens}[title={What you've built}]
One more word for everyday life at A2.
\end{grammarlens}

\subsection*{Your turn}
\begin{itemize}
\raggedright
  \item \emph{Say it:} \emph{mec}
  \item \emph{Say it:} \emph{mec}, once more
  \item \emph{Say it:} \emph{gaīkāle amārī śāḷānī mec hatī} — Yesterday our school had a match.
  \item \emph{Recall:} say the Gujarati for to hit, then the Gujarati for to bounce, then say \emph{mec} again
\end{itemize}

\begin{tcolorbox}[breakable,colback=teal!4,colframe=teal!35!black,title={Before you move on}]
What does \gu{મેચ} mean? (\textquotedblleft{}Match\textquotedblright{}.) Say it once more.
\end{tcolorbox}
\section[\texorpdfstring{\gu{હરાવવું} (harāvvũ)}{હરાવવું (harāvvũ)}]{\gu{હરાવવું} (harāvvũ) — to defeat, to beat}

\label{lesson:GU-C259-haravvu}



\begin{quote}
Before the new one: say the Gujarati for to bounce, then the Gujarati for match.
\end{quote}

\subsection*{You'll want to know: \gu{હરાવવું}}
\textbf{\gu{હરાવવું}} — \emph{harāvvũ} — \textquotedblleft{}to defeat, to beat\textquotedblright{}.

\begin{quote}
\textbf{\gu{અમે} \gu{તેમને} \gu{મેચમાં} \gu{હરાવ્યા}.} — \emph{ame temne mecmā̃ harāvyā} — \textquotedblleft{}We beat them in the match.\textquotedblright{}
\end{quote}

The causative of \textbf{\gu{હારવું}}, \emph{hārvũ}, \textquotedblleft{}to lose\textquotedblright{}: to make someone lose.

\begin{cousinweb}
\textbf{\gu{હરાવવું}} (\emph{harāvvũ}) is built on the taught \emph{hārvũ}, \textquotedblleft{}to lose\textquotedblright{}, with the causative \emph{-āv-}. CDIAL derives \emph{hārvũ} from Sanskrit \emph{hārayati}, \textquotedblleft{}has taken from oneself, loses\textquotedblright{} (CDIAL 14061).
\end{cousinweb}

\begin{grammarlens}[title={What you've built}]
One more word for everyday life at A2.
\end{grammarlens}

\subsection*{Your turn}
\begin{itemize}
\raggedright
  \item \emph{Say it:} \emph{harāvvũ}
  \item \emph{Say it:} \emph{harāvvũ}, once more
  \item \emph{Say it:} \emph{ame temne mecmā̃ harāvyā} — We beat them in the match.
  \item \emph{Recall:} say the Gujarati for to bounce, then the Gujarati for match, then say \emph{harāvvũ} again
\end{itemize}

\begin{tcolorbox}[breakable,colback=teal!4,colframe=teal!35!black,title={Before you move on}]
What does \gu{હરાવવું} mean? (\textquotedblleft{}To defeat, to beat\textquotedblright{}.) Say it once more.
\end{tcolorbox}
\section[\texorpdfstring{\gu{દોડવું} (doḍvũ)}{દોડવું (doḍvũ)}]{\gu{દોડવું} (doḍvũ) — to run}

\label{lesson:GU-C259-dodvu}



\begin{quote}
Before the new one: say the Gujarati for match, then the Gujarati for to defeat.
\end{quote}

\subsection*{You'll want to know: \gu{દોડવું}}
\textbf{\gu{દોડવું}} — \emph{doḍvũ} — \textquotedblleft{}to run\textquotedblright{}.

\begin{quote}
\textbf{\gu{હું} \gu{સવારે} \gu{મેદાનમાં} \gu{દોડ્યો}.} — \emph{hũ savāre medānmā̃ doḍyo} — \textquotedblleft{}I ran in the field in the morning.\textquotedblright{}
\end{quote}

\begin{cousinweb}
\textbf{\gu{દોડવું}} (\emph{doḍvũ}) is a tadbhava. It comes from Sanskrit \emph{drávati}, \textquotedblleft{}runs\textquotedblright{}, extended with \emph{-ḍ-}; CDIAL lists Gujarati \emph{doṛvũ}, \textquotedblleft{}to run\textquotedblright{} (CDIAL 6624).
\end{cousinweb}

\begin{grammarlens}[title={What you've built}]
One more word for everyday life at A2.
\end{grammarlens}

\subsection*{Your turn}
\begin{itemize}
\raggedright
  \item \emph{Say it:} \emph{doḍvũ}
  \item \emph{Say it:} \emph{doḍvũ}, once more
  \item \emph{Say it:} \emph{hũ savāre medānmā̃ doḍyo} — I ran in the field in the morning.
  \item \emph{Recall:} say the Gujarati for match, then the Gujarati for to defeat, then say \emph{doḍvũ} again
\end{itemize}

\begin{tcolorbox}[breakable,colback=teal!4,colframe=teal!35!black,title={Before you move on}]
What does \gu{દોડવું} mean? (\textquotedblleft{}To run\textquotedblright{}.) Say it once more.
\end{tcolorbox}
\section[\texorpdfstring{\gu{કસરત} (kasrat)}{કસરત (kasrat)}]{\gu{કસરત} (kasrat) — exercise, workout}

\label{lesson:GU-C259-kasrat}



\begin{quote}
Before the new one: say the Gujarati for to defeat, then the Gujarati for to run.
\end{quote}

\subsection*{You'll want to know: \gu{કસરત}}
\textbf{\gu{કસરત}} — \emph{kasrat} — \textquotedblleft{}exercise, workout\textquotedblright{}.

\begin{quote}
\textbf{\gu{મેં} \gu{સવારે} \gu{કસરત} \gu{કરી}.} — \emph{mẽ savāre kasrat karī} — \textquotedblleft{}I exercised in the morning.\textquotedblright{}
\end{quote}

\begin{cousinweb}
The origin of \textbf{\gu{કસરત}} (\emph{kasrat}) is not settled. It is usually linked with a Persian word of the same shape, but CDIAL has no entry for it and our sources do not confirm the link.
\end{cousinweb}

\begin{grammarlens}[title={What you've built}]
One more word for everyday life at A2.
\end{grammarlens}

\subsection*{Your turn}
\begin{itemize}
\raggedright
  \item \emph{Say it:} \emph{kasrat}
  \item \emph{Say it:} \emph{kasrat}, once more
  \item \emph{Say it:} \emph{mẽ savāre kasrat karī} — I exercised in the morning.
  \item \emph{Recall:} say the Gujarati for to defeat, then the Gujarati for to run, then say \emph{kasrat} again
\end{itemize}

\begin{tcolorbox}[breakable,colback=teal!4,colframe=teal!35!black,title={Before you move on}]
What does \gu{કસરત} mean? (\textquotedblleft{}Exercise, workout\textquotedblright{}.) Say it once more.
\end{tcolorbox}
\section[\texorpdfstring{\gu{કુસ્તી} (kustī)}{કુસ્તી (kustī)}]{\gu{કુસ્તી} (kustī) — wrestling}

\label{lesson:GU-C259-kusti}



\begin{quote}
Before the new one: say the Gujarati for to run, then the Gujarati for exercise.
\end{quote}

\subsection*{You'll want to know: \gu{કુસ્તી}}
\textbf{\gu{કુસ્તી}} — \emph{kustī} — \textquotedblleft{}wrestling\textquotedblright{}.

\begin{quote}
\textbf{\gu{મારા} \gu{ભાઈએ} \gu{ગામમાં} \gu{કુસ્તી} \gu{જોઈ}.} — \emph{mārā bhāīe gāmmā̃ kustī joī} — \textquotedblleft{}My brother watched wrestling in the village.\textquotedblright{}
\end{quote}

\begin{cousinweb}
\textbf{\gu{કુસ્તી}} (\emph{kustī}) is a loan, usually traced to Persian \emph{kushtī}, \textquotedblleft{}wrestling\textquotedblright{}. CDIAL has no entry for it.
\end{cousinweb}

\begin{grammarlens}[title={What you've built}]
One more word for everyday life at A2.
\end{grammarlens}

\subsection*{Your turn}
\begin{itemize}
\raggedright
  \item \emph{Say it:} \emph{kustī}
  \item \emph{Say it:} \emph{kustī}, once more
  \item \emph{Say it:} \emph{mārā bhāīe gāmmā̃ kustī joī} — My brother watched wrestling in the village.
  \item \emph{Recall:} say the Gujarati for to run, then the Gujarati for exercise, then say \emph{kustī} again
\end{itemize}

\begin{tcolorbox}[breakable,colback=teal!4,colframe=teal!35!black,title={Before you move on}]
What does \gu{કુસ્તી} mean? (\textquotedblleft{}Wrestling\textquotedblright{}.) Say it once more.
\end{tcolorbox}
